% AI-written, AI cross-reviewed; no human sign-off; not a claim that RH is proved
% Atlas unit D479. Local manuscript v3, item G.499, heading on printed PDF p. 281.
% GlobalTheory.Top.link_gl
% paper2/decomposition_tables.tex:4136–4141. Verbatim excerpt; SHA-256 46dffa67cc0eda5e9297d534fa1f25bb0a98bde3d5b5918d381b82d5535ea7ec.
% Original TeX uses the macros and cross-reference keys of the source manuscript.

\begin{formalstatement}[D479]\label{decomp:D479}
On the enriched class D3(3,A), duty, tied real-place second-layer data, the mass-ratio estimate and the genus-mass hypotheses give two uniform implications. Chosen data of exponent 3 give fullness and the prescribed degree/rank relation. Tuple and discriminant-size bounds give the global comparison with constants 2, CH, 3, 2, 0.
\end{formalstatement}
{\small\raggedright\noindent\textsf{\bfseries Project-specific}\quad \nolinkurl{GlobalTheory.Top.link_gl}\par\noindent \nolinkurl{Development/AbelianVarietyTheory/Proofs/GlobalInterpolation/LinkGl.lean:14--31}\par\noindent SHA-256 \nolinkurl{50ea47933db45c6a1bd8444f061acfb36555e58406fce4b4a74fdf5df0575ae2}\par\noindent Dominated 11,198\quad Disjoint 11,198\quad terminal.\par}
The mass comparison must control the chosen family, its ranks and the same constants in the global lattice estimate. The unitary-lattice source does not assert these implications.
\hyperref[lit:S055]{Published sources and search S055}.
